\section{¿Cómo Debemos Investigarlo?}
De acuerdo con los criterios revisados en clase, el tipo de investigación se
elige a partir del problema y de la pregunta planteados en las secciones
anteriores. Cada decisión se declara y se justifica enseguida.

\subsection{Finalidad}
La investigación es \textbf{aplicada}. Se entiende por aplicada la
investigación dirigida primordialmente a un fin práctico específico, en
contraste con la investigación pura o básica, cuyo propósito es adquirir
conocimiento sin un objetivo de aplicación predeterminado
\parencite{oecd2015frascati}; su valor está en producir evidencia para
resolver problemas concretos de la práctica \parencite{vargascordero2009aplicada}.
Este proyecto se ubica en el primer caso: no se pretende generar conocimiento
teórico nuevo sobre los códigos QR ni sobre la teoría de colas, temas que ya
cuentan con literatura consolidada, sino emplear ese conocimiento existente
para resolver un problema situado, el control de acceso de un plantel
determinado. El propósito del estudio es producir una decisión práctica —si
conviene o no adoptar el sistema en el resto de los accesos— y un producto
utilizable, el prototipo del Apéndice~D. Sus resultados valen para el caso
estudiado y no se presentan como una generalización.

\subsection{Enfoque}
El enfoque es \textbf{mixto, con predominancia cuantitativa}
\parencite{hernandezsampieri2014metodologia}, porque el problema exige dos
clases distintas de información. Los indicadores de tiempo
de validación, longitud de fila e incidencias de identidad (encuesta del
Apéndice~A y hoja de observación del Apéndice~C) son magnitudes medibles y se
analizan de forma cuantitativa, pues la pregunta de investigación pide saber
\emph{en qué medida} cambian. La entrevista del Apéndice~B aporta información
cualitativa que ninguna medición recoge: las razones, los criterios y las
dificultades que el personal de vigilancia y de Control Escolar expresa con
sus propias palabras.

\subsection{Alcance}
El alcance es \textbf{descriptivo y explicativo}. Es descriptivo en su primera
etapa, de carácter diagnóstico, porque el nivel de conocimiento disponible
sobre el acceso del TESCHA es nulo y hace falta documentar la situación actual
—cuánto se espera, cuántos entran, con qué frecuencia falla la
identificación— para establecer una línea base. Es explicativo en su segunda
etapa, porque una vez instalado el prototipo se busca establecer si el cambio
de mecanismo de validación explica la reducción observada en el tiempo de
ingreso y en las incidencias de identidad, y no solo constatar que ocurrió.

\subsection{Obtención de Información}
La obtención de información es \textbf{mixta}, documental y de campo. Es
documental en la fundamentación: la elección del código QR frente a NFC y al
reconocimiento facial, el esquema de firma HMAC y el diseño cuasiexperimental
se sustentan en fuentes ya publicadas, revisadas en las secciones anteriores.
Es de campo en la parte central del estudio, porque los datos que dan
respuesta a la pregunta de investigación no existen en ninguna fuente
documental: no hay registro alguno del flujo de acceso del TESCHA, de modo
que deben levantarse en el sitio, en el acceso piloto y en el horario en que
el problema se presenta.

\subsection{Técnica e Instrumento de Recolección de Datos}
Se emplean tres pares de técnica e instrumento, cada uno elegido por el tipo
de información que se necesita obtener:
\begin{itemize}
	\item \textbf{Encuesta}, con un \textbf{cuestionario} de preguntas cerradas
	      y escala de Likert (Apéndice~A). Se elige porque la percepción de espera, la
	      frecuencia de olvido de credencial, la disposición a usar el sistema y la
	      viabilidad tecnológica del estudiante son datos que solo el propio
	      estudiante puede reportar, y se requieren de una población numerosa en poco
	      tiempo, con respuestas comparables entre sí y de forma anónima.
	\item \textbf{Entrevista semiestructurada}, con una \textbf{guía de
		      entrevista} (Apéndice~B). Se elige porque el personal de vigilancia y de
	      Control Escolar son informantes clave, pocos en número, que poseen
	      información que ningún cuestionario anticipa: el procedimiento real de
	      revisión, los casos excepcionales y las condiciones institucionales para
	      tratar datos de los estudiantes. El formato semiestructurado permite
	      repreguntar sobre lo que aparezca.
	\item \textbf{Observación estructurada no participante}, con una \textbf{hoja
		      de registro} (Apéndice~C). Se elige porque el tiempo de revisión y la
	      longitud de la fila son hechos observables que deben medirse
	      directamente con cronómetro y conteo; preguntarlos produciría estimaciones
	      subjetivas, no mediciones. Es además el único instrumento que se aplica dos
	      veces con el mismo formato, lo que hace comparables la línea base y el
	      piloto.
\end{itemize}

\subsection{Control de Variables}
En cuanto al control de variables, la investigación es \textbf{experimental},
aunque no en sentido estricto, sino \textbf{cuasiexperimental}. Sí hay
manipulación deliberada de la variable independiente —el mecanismo de
validación de identidad, que pasa de la revisión visual de credenciales al
escaneo del QR dinámico— y se observa su efecto sobre las variables
dependientes: el tiempo promedio de validación por estudiante, la longitud
máxima de la fila y el número de incidencias de identidad. Lo que no es
posible es el control propio de un experimento puro: no hay asignación
aleatoria de los estudiantes ni grupo de control, porque el acceso piloto
atiende a quien llega. Se mantienen constantes las condiciones que sí están al
alcance del equipo —el mismo acceso, el mismo turno, la misma franja horaria,
los mismos observadores y el mismo instrumento en ambas mediciones— y se
registra el número de puntos de revisión activos ($c$) para no atribuir al
prototipo una mejora que provenga de haber abierto más filas.

\subsection{Temporalidad}
Según el criterio temporal, que distingue entre recolectar datos en un solo
momento y recolectarlos en varios puntos en el tiempo para hacer inferencias
sobre el cambio \parencite{hernandezsampieri2014metodologia}, esta
investigación es \textbf{longitudinal}, con mediciones repetidas. El fenómeno
no se mide una sola vez, sino en dos series separadas por la intervención:
cinco días hábiles consecutivos de observación antes del piloto y una segunda
serie equivalente durante él, sobre el mismo acceso y con la misma hoja de
registro. Esta temporalidad la exige la pregunta de investigación, que pide
comparar un antes y un después; un estudio transversal, de una sola medición,
permitiría describir la situación actual pero no atribuir ningún cambio al
prototipo.

Conviene precisar el uso del término, porque el eje
transversal/longitudinal suele desarrollarse para los diseños no
experimentales, en los que el investigador observa sin administrar ningún
estímulo, y sus subtipos —como el diseño de panel— se definen en ese marco
\parencite{hernandezsampieri2014metodologia}. Aquí sí hay un estímulo
administrado por el equipo, el prototipo, de modo que el diseño pertenece a la
familia cuasiexperimental descrita más adelante. Por eso «longitudinal» se
emplea en este trabajo en su sentido temporal —varias mediciones a lo largo
del tiempo sobre las mismas unidades de observación— y no como la etiqueta de
un diseño no experimental de panel.

\subsection{Diseño}
En consecuencia, el diseño es de campo y cuasiexperimental, del tipo
preprueba/posprueba con un solo grupo —el acceso piloto— y sin grupo de
control \parencite{campbellstanley1963, hernandezsampieri2014metodologia},
dado que no es posible operar dos accesos idénticos en paralelo bajo las
mismas condiciones. Un diseño equivalente se ha empleado recientemente para
evaluar, también sin grupo de control, un piloto de intervención basado en
código QR dentro de una institución \parencite{ortiz2026qrpilot}. Entre las
amenazas a la validez interna de este diseño destacan la historia y la
maduración \parencite{cranmer2017pretestposttest} —por ejemplo, un cambio de
horario escolar o la simple costumbre de los estudiantes podrían alterar los
tiempos sin que el prototipo tenga parte en ello—, por lo que los resultados
se interpretarán como evidencia orientativa y no concluyente. Los tres
instrumentos se aplican antes del piloto para obtener la línea base y de
nuevo durante él para comparar los resultados, como se detalla en el Plan de
Acción Preliminar (Apéndice~D).
